\noindent\textbf{Elección del procedimiento.} La ley de ondas largas dada es $c=\sqrt{gH}$. El tiempo de viaje es longitud dividida por esa celeridad; solo después se convierte de segundos a horas. Se mantienen las hipótesis de profundidad uniforme y longitud de onda grande frente a $H$.

\begin{enumerate}

\item Calcular primero la velocidad de propagación.

\[\fitmath{c=\sqrt{g\cdot H}}\]

\[\fitmath{c=\sqrt{(9.81)\cdot (4000)}\approx 198.0909\,\mathrm{m/s}}\]

\item Convertir 2000 km a 2000000 m.

\[\fitmath{\mathrm{ts}=\frac{\mathrm{distance}}{c}}\]

\[\fitmath{\mathrm{ts}=\frac{(2000000)}{(198.0908882)}\approx 10096.38\,\mathrm{s}}\]

\item Un resultado en segundos se divide por 3600 para expresarlo en horas.

\[\fitmath{\mathrm{th}=\frac{\mathrm{ts}}{3600}}\]

\[\fitmath{\mathrm{th}=\frac{(10096.37555)}{3600}\approx 2.804549\,\mathrm{h}}\]

\end{enumerate}
